\subsection{Problem Formulation and Design Requirements}
\label{sec:control-problem}

Equation~\ref{eq:draft-budget} exposes two distinct control needs.
Remaining-sequence admission controls current demand by
selecting which optional effects contribute to $T_i^{\mathrm{effects}}$, while
preserving the required and mandatory tail. It must keep
$T_i^{\mathrm{effects}}+T_i^{\mathrm{tail}}\leq b_i$. Admission cannot,
however, restore budget already consumed before the Draft Sequence obtains the
worker.

If captures arrive faster than the single worker completes them, repeated
admission decisions alone cannot stop $b_i$ from shrinking across the
burst. The framework must also control when the next capture may begin.
Delaying the application-facing \texttt{captureAvailable} callback does not
extend the current deadline; it lets queued Draft work progress before the next
capture---and its timeout budget---starts.

The industrial objective is consequently lexicographic: (1) prevent every
Capture Timeout, because deadline misses are release blocking; (2) retain the
preferred Draft effects on as many captures as the runtime state safely allows;
and (3) minimize added callback delay and its effect on shot-to-shot
responsiveness. Section~\ref{sec:approach} presents the shared tail model and
the coordinated admission and pacing decisions that address these requirements.
